\vspace{2\baselineskip}
\section{Conclusion}

\section{Companies that systematically lose public tenders in which they participate are not expected in a competitive environment. An analysis of the relationship between performance variables and the participation of these companies in public bids in São Paulo State between 2009 and 2019 indicates that the negotiated price was 10% higher; however, the number of participants was 32% higher, and the number of bids was 29% higher.}

\section{This apparent paradox, in which one variable indicates less competitive pressure (i.e., negotiated prices) while others indicate more competition (i.e., number of participants and number of bids), is consistent with the behavior of a cartel that seeks to avoid detection. This scenario is the function that would intuitively be expected of a frequent loser, that is, to simulate a higher level of competition, drive away competitors, and avoid detection by control agencies. Precisely because of this characteristic, frequent losers can be used as markers in a cartel screening method.}

\section{This proposal presents some virtues common to other screening methods and addresses two limitations common to most available models. It is a simple application method that requires only data already available to competition authorities and control agencies.}

\section{The identification of frequent losers may occur before the public tender on the occasion of the first stage of the bid, in which the participants are defined. Unlike all available screening methods based on prices and quantities observed after the cartel's materialization, frequent losers may signal suspicious behavior before bidding occurs, reducing the costs associated with a defrauded event.}

\section{This property stems from a characteristic of cartels in public tenders: the agreement among bidders occurs before the tender takes place, and it is possible to observe elements that may signal fraud before the conclusion of the bid. The presence of frequent losers also emphasizes the separation between explicit collusion and tacit collusion since the behavior identified is consistent with the deliberate act of concealing competition.}

\section{There are, however, some limitations in the use of these screening methods; some of these limitations are remedial, and another is common to all. The proposed method identifies suspicious conduct only in cases of bidding cartels, where some mechanism for transferring the benefits of the cartel among its members is possible. This is still a relevant subset of cases, but it does not apply to all cartels. However, this limitation is remedial because the method can be associated with any other screening method and can act in a complementary manner, as suggested by .}

\section{common problem with all screening methods is that once they become known to cartel participants, the cartel participants will modify their behavior to avoid detection (Schinkel, 2014). However, this is not a reason to rule out screening methods; on the contrary, it is necessary to develop them continuously, and the cumulative process of screening methods is associated with increasing costs for cartel participants to avoid detection. Eventually, this cumulative process may result in the deterrence of such an antitrust offense.}
